\section{The determinant comparison}\label{sec:comparison}
We recall the determinant and the estimates from~\cite{chung498} that will
be used. The comparison applies to every nonzero maximal minor. Thus the
rank argument in the next section need only establish that such a minor
exists; it need not identify its columns.

\subsection{The matrix and its normalization}
Assume that $h(-D)=1$ and $D>331$. Write
$\omega=(1+\sqrt{-D})/2$ and $\OO_D=\mathbb Z[\omega]$.
Choose integers $K\geq2$, $S\geq1$ and $1\leq J\leq R:=KS$, and put
\begin{equation}\label{eq:triangle}
 \Lambda=\{(r,s)\in\mathbb Z_{\geq0}^2:r+Ks<R\},\qquad
 M=|\Lambda|=\frac{R(S+1)}2.
\end{equation}
Let $z_1,\ldots,z_n$ be the $n$ elements of $\OO_D$ nearest to the origin,
with ties broken arbitrarily. We use the Hasse derivatives in $z$, with the
second variable $w$ held fixed: $\partial_z^{[j]}P$ is the coefficient of
$t^j$ in $P(z+t,w)$. The matrix is
\begin{equation}\label{eq:jet-matrix}
 V_{(r,s),(i,j)}=\binom rjz_i^{r-j}\bar z_i^{\,s},
 \qquad (r,s)\in\Lambda,\quad 1\leq i\leq n,\quad 0\leq j<J.
\end{equation}
The entry is zero if $j>r$. Dividing ordinary derivatives by $j!$ removes their factorial
factors while retaining algebraic-integer entries. For any selection
$C$ of $M$ columns, let $\Delta_C=\det V_C$.

The nearest-node estimate of~\cite[Lemma 2.5]{chung498} gives
\begin{equation}\label{eq:node-disk}
 |z_i|\leq\rho_i:=\sqrt{\frac{i\sqrt D}{2\pi}}+\delta,
 \qquad \delta=\frac{\sqrt{D+4}}4.
\end{equation}
The lattice has covolume $\sqrt D/2$ and covering radius
at most $\delta$: its horizontal rows have vertical spacing $\sqrt D/2$,
and the points on each row have spacing one. The disk of radius
$\rho-\delta$ is therefore covered by Voronoi cells of lattice points in
the disk of radius $\rho$. Comparing areas shows that this latter disk
contains at least $2\pi(\rho-\delta)^2/\sqrt D$ lattice points, proving
\eqref{eq:node-disk}.

Put $\rho=\rho_n$, $\theta=M/(nJ)$, and
\begin{equation}\label{eq:normalization}
 W=\frac{R^2S}{3},\qquad W_1=2\sum_\Lambda r,\qquad W_2=2\sum_\Lambda s.
\end{equation}
The ratio $\theta$ is the proportion of available columns used by a maximal
minor. In particular, $\theta\leq1$ is necessary for full row rank. Direct
summation over the triangle gives
\begin{equation}\label{eq:radius-coefficient}
 \frac{W_1+W_2}{W}
 =1+\frac1K+\frac{3(K-1)}{2KS}+\frac{K-5}{2KS^2}.
\end{equation}

\subsection{The arithmetic estimate}\label{subsec:primes-import}
We use the local divisibility formulas of the earlier paper, with all
quantities evaluated at the present parameters. Their general derivation is
given in Appendices~A and~B of~\cite{chung498}.
We record the hypotheses and the meaning of the resulting terms, since the
finite comparisons in that paper use different parameters.

Suppose throughout this subsection that $D>163$ is prime,
$D\equiv3\pmod8$, and $h(-D)=1$. Put $m_D=(D+1)/4$.
Let $P$ be the prime cutoff in the arithmetic sums. The parameters used
here satisfy
\begin{equation}\label{eq:arithmetic-hypotheses}
 \begin{gathered}
 2\leq K<D/4,\quad R=KS,\quad 1\leq J\leq R,\quad M\leq nJ,\\
 P\geq\max(S,J-1),\qquad
 F_{\max}:=\lfloor P/m_D\rfloor<m_D^4.
 \end{gathered}
\end{equation}
The local divisibility estimates give a number $\arith$ and a derivative
loss $b(j)$ such that every nonzero maximal minor satisfies
\begin{equation}\label{eq:arithmetic-interface}
 W\arith-2\sum_C b(j)\leq\log|\det V_C|^2.
\end{equation}
We recall their construction to specify how the estimates apply at the
present parameters. An element is \emph{primitive} if it is not divisible
in $\OO_D$ by a rational integer greater than one.

All primes below $K$ are then inert. The small-prime estimates apply without
exception, and the cutoff $P$ includes the derivative debits, the ramified
prime, and the primes contributing to the conjugate-variable estimates.
Summing the local row bounds gives
\begin{equation}\label{eq:arithmetic-parts}
 \arith=\arith_0-L_{\rm norm}+C_2+C_3.
\end{equation}
Here $\arith_0$ provisionally treats primes up to $P$ as inert.
Class number one represents every non-inert prime as the norm of an
algebraic integer. Summing over primitive norm representatives therefore
gives an upper bound $L_{\rm norm}$ for the necessary correction, retaining
the split and ramified local contributions. This sum also includes composite
norm values. The nonnegative terms $C_2,C_3$ correct this overestimate for
products of two or three prime factors, with the representation
multiplicities supplied by unique factorization. Omitting other composite
corrections only weakens the lower bound.

Some split-prime column contributions are retained individually. Let
$\mathcal P_s(D)$ be the split primes at most $q_{\rm cap}$, where the
cutoff satisfies
\begin{equation}\label{eq:slope-cutoff}
 q_{\rm cap}\leq\min(J-1,F_{\max}m_D,m_D^4-1).
\end{equation}
Define the remaining derivative loss and its factorial remainder by
\begin{equation}\label{eq:arithmetic-debit}
 \begin{split}
 b(j)&=\sum_{q\leq j}\lfloor j/q\rfloor\log q
       -\sum_{q\in\mathcal P_s(D)}\lfloor j/q\rfloor\log q,\\
 F(j)&=\log(j!)-b(j)
       =\sum_q\sum_{a\geq2}\lfloor j/q^a\rfloor\log q
        +\sum_{q\in\mathcal P_s(D)}\lfloor j/q\rfloor\log q.
 \end{split}
\end{equation}
Sums indexed by $q$ are over primes. The split-prime contributions outside this
window are bounded uniformly over column selections and already included
in \eqref{eq:arithmetic-parts}.

The prime sums in the window are bounded using norm representatives and
composite products. Every prime factor of a primitive norm is at least
$m_D$. Thus a composite primitive norm in \eqref{eq:slope-cutoff} has at most
three prime factors, counted with multiplicity, and each factor is at
most $F_{\max}$. To obtain the factor list up to $F_{\max}$, the finite enumeration
first removes reducible primitive elements. A reducible
primitive element of norm at most $F_{\max}<m_D^4$ has an irreducible factor
of norm at most $\sqrt{F_{\max}}<m_D^2$, and a remaining factor of norm at
most $F_{\max}/m_D$. Appendix~B of~\cite{chung498} proves that these ranges construct the complete factor list. Consequently the subtraction includes every
composite in the window, even for $P>m_D^4$.

Using Appendix~B of~\cite{chung498}, we choose a nonnegative lower bound $\sigma$ for
\[
 \sum_{\substack{q,\ a\geq2\\q^a\leq J-1}}\frac{\log q}{q^a}
       +\sum_{q\in\mathcal P_s(D)}\frac{\log q}{q},
\]
and upper bounds $h_j$ for $H(j)=j\sigma-F(j)$.
These bounds determine $c$ by the capacity calculation below.
The explicit finite weights, slope bounds and remainder estimates are
those in Appendix~B of~\cite{chung498}. The hypotheses above hold for all three parameter
families below.


\subsection{Combining the derivative terms}
Legendre's formula writes $\log(j!)$ as
$\sum_q\sum_{a\geq1}\lfloor j/q^a\rfloor\log q$. Its first-power
terms cancel the corresponding derivative loss in $b(j)$, leaving the
higher prime powers and split-prime contributions in $F(j)$. We use the
remainder $H(j)=j\sigma-F(j)$ defined above,
and choose a constant $c$ satisfying
\begin{equation}\label{eq:column-remainder}
 2\sum_{(i,j)\in C}H(j)\leq Mc\qquad(|C|=M).
\end{equation}
This bound uses the fact that each derivative order is available at only
$n$ nodes. Specifically, if $h_{(1)}\geq\cdots\geq h_{(J)}$ are upper
bounds for $H(0),\ldots,H(J-1)$ arranged in decreasing order, and
$M=nt+r$ with $0\leq r<n$, one may take
\[
 Mc=2n\sum_{a=1}^{t}h_{(a)}+2r h_{(t+1)}.
\]
The last term is omitted when $r=0$. The formulas in Appendix~B of~\cite{chung498} construct $\sigma$ and these upper bounds from prime powers and norm
values, so the calculation does not require identifying the split primes
at the candidate discriminant.

Set $u_i=z_i/\rho$ and $\ell=\log(R/\rho)-\sigma$. Define the normalized
matrix
\begin{equation}\label{eq:residual-matrix}
 B_{(r,s),(i,j)}(\ell)
 =\frac{\fall rj}{R^j}e^{\ell j}u_i^{r-j}\bar u_i^{\,s},
\end{equation}
where $\fall rj=r(r-1)\cdots(r-j+1)$. Exact diagonal scaling gives
\[
 B_C(\ell)=\operatorname{diag}_{\Lambda}(\rho^{-r-s})\,
 V_C\,\operatorname{diag}_C(j!e^{-\sigma j}).
\]
Taking squared absolute determinants and using \eqref{eq:arithmetic-interface} and
\eqref{eq:column-remainder}, we obtain
\begin{equation}\label{eq:comparison-before-size}
 W\arith\leq (W_1+W_2)\log\rho+
                \log|\det B_C(\ell)|^2+Mc.
\end{equation}
Since $e^{\ell j}$ scales column $(i,j)$ and $j\geq0$, replacing
$\ell$ by an upper bound only increases $|\det B_C(\ell)|^2$. It
therefore suffices to bound the normalized determinant uniformly over $C$.

\subsection{The imported size estimate}
We use the Meixner and Schur estimates from~\cite[Section~4]{chung498}.
Their role is to retain the correlations between derivative columns at
one node. We recall the quantities entering the bound so that its
application at the new parameters is explicit.

Choose $a>0$, $0<v<1$, $\beta<0$, and put $b=aR$. Weight row $(r,s)$ by
$e^{-\lambda_{r,s}/2}$, where
\[
 e^{-\lambda_{r,s}}=\frac{(b)_r}{r!}v^r e^{-\beta s},
 \qquad (b)_r=b(b+1)\cdots(b+r-1).
\]
Let $Y=\operatorname{diag}(e^{-\lambda_{r,s}/2})B(\ell)$, and let $G_i=Y_i^*Y_i$
be the Gram matrix of its $J$ columns at node $i$. Fischer's inequality gives
\begin{equation}\label{eq:fischer}
 \log|\det B_C(\ell)|^2
 \leq\sum_\Lambda\lambda_{r,s}
       +\sum_i\log\det G_i[C_i,C_i],
\end{equation}
where $C_i$ is the set of orders selected at that node. By~\cite[Lemma~4.1]{chung498}, each principal minor
of $G_i$ is a polynomial with nonnegative coefficients in $q=|u_i|^2$;
therefore it increases when the node is moved radially outward.

Write $x_+=\max(0,x)$. At an upper squared radius $q>0$, set
\[
 d_q=(\log q-\beta)_+,\qquad v_q=ve^{-d_q/K},\qquad
 w=v_qq,\qquad \mathcal B_q=S e^{Sd_q}.
\]
After a diagonal unitary conjugation removing the node angle, the Gram
matrix is bounded in the positive semidefinite order by
\begin{equation}\label{eq:moment-matrix}
 \mathcal G_{jk}(q)=\mathcal B_q
 \left(\frac{e^\ell}{R\sqrt q}\right)^{j+k}
 \sum_{r\geq0}\frac{(b)_r}{r!}w^r\fall rj\fall rk,
 \qquad 0\leq j,k<J.
\end{equation}
The conjugation preserves every principal minor. Here $0<w<1$ for the
parameters below. The negative-binomial moments in
this formula have an explicit factorization in Meixner polynomials. This
supplies two bounds for an arbitrary principal minor: one uses the largest
singular values of an affine composition operator, and the other its
complementary singular product. A third bound sums all principal minors of
the prescribed size, using the Schur-function Cauchy identity. These are
respectively equations (4.10), (4.11), and (4.12) of~\cite{chung498}.
Each holds for every subset of derivative orders; no particular subset is
assumed to maximize the determinant.

Let $F_q(m)$ be the minimum of these three upper bounds for the logarithm
of an $m$-column principal minor, allowing the finite parameter choices in
Appendix~C of~\cite{chung498}. For a chosen $a_0>0$, take
$\Phi(q)\geq0$ such that
\begin{equation}\label{eq:H-occupancy}
 F_q(m)\leq a_0Rm+R^2\Phi(q)\qquad(0\leq m\leq J).
\end{equation}
We divide the range of $m/R$ into $N_m$ equal intervals and establish
the bound on each whole interval before minimizing over estimates.
For the singular-value bounds, the sum of the $m$ largest column prices
is bounded using $N_z$ equal intervals of $j/R\in[0,\varepsilon]$,
where $\varepsilon=J/R$. Appendix~C of~\cite{chung498} bounds the price and the number of orders in each
interval. Summing \eqref{eq:H-occupancy} in
\eqref{eq:fischer} uses $\sum_i|C_i|=M$, rather than allowing every node to
choose its worst number of columns independently.

To bound the radii, put
\[
 \eta=\frac{\delta}{\sqrt{n\sqrt D/(2\pi)}},\qquad
 Q(t)=\left(\frac{\sqrt t+\eta}{1+\eta}\right)^2,
 \qquad \varepsilon=J/R.
\]
Then $|u_i|^2\leq Q(i/n)$. Divide $(0,1]$ into $N_t$ equal intervals and
move the nodes in the $k$th interval to squared radius $Q(k/N_t)$.
Each interval contains at most $n/N_t+1$ nodes. The resulting bound is
\begin{equation}\label{eq:metric-mesh}
 U_{\rm var}=U_{\rm row}+\frac{3a_0(S+1)}{2S}
 +3\left(\frac{1+1/S}{2\theta\varepsilon N_t}+\frac1S\right)
       \sum_{k=1}^{N_t}\Phi(Q(k/N_t)),
\end{equation}
where $U_{\rm row}$ is an upper bound for
$W^{-1}\sum_\Lambda\lambda_{r,s}$. In particular,
$W^{-1}\log|\det B_C(\ell)|^2\leq U_{\rm var}$ for every $C$.
For the finite formulas in Appendix~C of~\cite{chung498},
we require $R>3$, $aR>1$, $0<a\leq1$, $0<\varepsilon\leq1$,
$\varepsilon(a+2\varepsilon)\leq1$, and
\[
 -(a+1)\log(a+1)+a\log a-\log v<0.
\]
The last condition gives the signed row-sum bound. For the Schur bound at
radius $q$, put
\[
 \zeta=\ell+\tfrac12\log v_q-\log(1-w),\qquad z_*=a^{-2}e^{-2\zeta}.
\]
Its finite evaluation is used when $z_*<1$ and
$\varepsilon(a+\varepsilon)z_*<a^2(1-z_*)^2$. These are respectively the
moment and concavity conditions in Appendix~C of~\cite{chung498}. The other
two bounds remain available wherever the latter conditions do not hold.

\begin{proposition}[Determinant comparison]\label{prop:comparison}
Assume the arithmetic estimates \eqref{eq:arithmetic-interface} and
\eqref{eq:column-remainder}, and the size estimate \eqref{eq:metric-mesh}.
If $V$ has full row rank, then
\begin{equation}\label{eq:master}
 \arith\leq\metric:=\frac{W_1+W_2}{W}\log\rho+U_{\rm var}
                       +\frac{Mc+2M\log M}{W}.
\end{equation}
\end{proposition}
\begin{proof}
Full row rank supplies a nonzero minor. Apply
\eqref{eq:comparison-before-size} and the uniform size bound to that minor.
The nonnegative term $2M\log M/W$ retains the earlier numerical
normalization~\cite{chung498}.
\end{proof}
Thus a positive lower bound for $G:=\arith-\metric$ excludes the proposed
field. The arithmetic and size estimates depend on the actual node count
$n$ throughout. In particular, reducing $n$ changes both the disk radius
and the capacity bound for the derivative orders.
